\scanpage{175}\mspage{30}%
phraser une intuition topologique extrêmement simple, que je vais expliciter, en supposant pour fixer les idées qu'on est sur le corps de base $\mathbb C$ et qu'on en profite pour travailler avec
% Compact isolated point at whole-image (171,423), left of L5; origin and role uncertain.
\raisebox{0pt}[0pt][0pt]{\parbox{0pt}{\figure{\begin{tikzpicture}[overlay,x=1pt,y=1pt]\fill (-7,0) ellipse[x radius=.35pt,y radius=.5pt];\end{tikzpicture}}}}%
les espaces loc. compacts habituels. L'idée topologique derrière l'hypothèse d'équisingularité, c'est qu'il existe un \uline{voisinage tubulaire $T$} de $X_0$ dans $X$ se rétractant sur $X_0$, et tel que $(X_0,T)$ au dessus de $X_0$ soit un fibré \unsure{loclement|localement} trivial, donc $T-X_0$ au dessus de $X_0$ est loc trivial. En fait, \deleted{\unsure{on pourra con|on pourra cons}} \inserted{si $\partial T$ est le} ``bord'' de $T$, qui sera aussi fibré loc. trivial sur $X_0$, $T$ sur $X$ est le fibré en cônes déduit ($\simeq(\partial T\times I)\amalg_{\partial T}X_0$, où $I=[0,1]$, $\partial T%
% Compact isolated point at whole-image (702,1183), above the first relation shaft; not a mathematical dot.
\raisebox{1.3ex}[0pt][0pt]{\parbox{0pt}{\figure{\begin{tikzpicture}[overlay,x=1pt,y=1pt]\fill (4,0) -- (4.7,.6) -- (4.8,-.5) -- cycle;\end{tikzpicture}}}}%
\mathrel{\to}\partial T\times I$ défini par $x\mathrel{\unsure{\mapsto|\vdash}}(x,1)$, et $\partial T\mathrel{\unsure{-|\to}}X_0$ est la projection) donc $\mathring T=T-X_0\simeq\partial T\times[0,1[$ est $X_0$-homotope à $\partial T$. Si $X_0$ et $X_1$ sont non singuliers, il en est de $\unsure{=|\overline m}$ de $\mathring T$ et $\partial T$, qui sont des fibrés topologiquement lisses sur $X_0$. D'ailleurs, $\deleted{\unsure{X_1|X}}$ posant $\widetilde X_1=X_1-\mathring T$, \deleted{$\unsure{\widetilde X_1|X_1}$ dev.} l'inclusion $\widetilde X_1%
% Short isolated dark tick at whole-image (259,1797), below this source object; origin and role uncertain.
\raisebox{-2.3ex}[0pt][0pt]{\parbox{0pt}{\figure{\begin{tikzpicture}[overlay,x=1pt,y=1pt]\draw[line width=.55pt] (0,.6) -- (.2,-.6);\end{tikzpicture}}}}%
\hookrightarrow X_1$ est une équivalence d'homotopie, et $X$ se réconstitue, à homéomorphisme près, à partir \unsure{du|des} diagramme d'espaces
\diagram{\begin{tikzcd}[ampersand replacement=\&,column sep=small,row sep=large]
(\mathring T=T-X_0\underline{\approx}) \& \partial T \arrow[r,"\unsure{J|j}","\text{inclusion}"'] \arrow[d,"p","\text{fibration}"'] \& \widetilde X_1\,(\underline{\approx}X_1)\\
\& X_0 \&
\end{tikzcd}}
comme la somme amalgamée. En termes de ce diagramme d'espaces; le foncteur $\varphi$ plus
% archive sheet number: 175
